% Table I — positioning, by what each route presupposes.
% Used by sections/02-related-work.
%
% NOTE ON FRAMING. This table merges two that used to be separate: a
% related-work grid (Work / Emphasis / Distinction) and a route comparison.
% They answered the same question — how is this different from what exists —
% for papers and for approaches respectively, and the page budget does not
% carry both. Rows are ROUTES to a control interface, each citing its
% representative work, because that is the axis on which they actually differ.
%
% A latency column was deliberately not included: rosbridge is the transport
% this work runs on, not a competitor, so comparing throughput would compare a
% thing with itself. What separates these routes is what they require the
% vendor to have provided first, which is the last column.
\begin{table}[!t]
\caption{Routes to a control interface on a closed service robot, and what each presupposes}
\label{tab:positioning}
\centering
\scriptsize
\setlength{\tabcolsep}{3pt}
\begin{tabular}{@{}>{\raggedright\arraybackslash}p{1.75cm}ccc>{\raggedright\arraybackslash}p{2.3cm}@{}}
\toprule
\textbf{Route} & \textbf{Offl.} & \textbf{Insp.} & \textbf{N.-d.} & \textbf{Presupposes} \\
\midrule
Vendor SDK~\cite{glauser2023interop} & varies & no & yes &
A published SDK, kept current by the vendor \\[2pt]
Vendor cloud API & no & no & yes &
A subscription and a network path off site \\[2pt]
Bridge integration~\cite{crick2012rosbridge,aragao2016yarp} & yes & yes & yes &
Control of the \ros{} graph, and known message definitions \\[2pt]
Re-architecting with vendor support~\cite{kim2005reengineering} & yes & yes & yes &
Source code and access to the vendor's engineers \\[2pt]
Companion-app recovery~\cite{giese2018xiaomi,muhe2010kuka} & --- & yes & yes &
An artifact that already speaks the protocol \\[2pt]
Firmware replacement & yes & yes & \textbf{no} &
Root access; voids warranty, unrecoverable if it fails \\[2pt]
\textbf{This work} & \textbf{yes} & \textbf{yes} & \textbf{yes} &
\textbf{Only lawful possession of the robot} \\
\bottomrule
\end{tabular}
\end{table}
